\documentclass[10pt,a4paper]{amsart}
\usepackage[margin=19mm]{geometry}
\usepackage{amsmath,amssymb,mathtools}
\usepackage[scr=rsfs,cal=euler]{mathalfa}
\usepackage{xfrac}
\usepackage[unicode,hidelinks,bookmarks=false]{hyperref}
\newcommand{\ie}{i.e.\ }
\newcommand{\cf}{cf.\ }
\newcommand{\resp}{respectively\ }
\newcommand{\Subsection}{\subsection}
\providecommand{\nobreakdash}{\nobreak\hbox{-}}
\setlength{\emergencystretch}{2em}
\multlinegap0pt
\usepackage{bm}
\usepackage{bbm}
\usepackage{dsfont}
\usepackage{xspace}
\usepackage{cancel}
\usepackage{mathtools}
\usepackage{amsthm}
\makeatletter
\@ifundefined{prop}{\newtheorem{prop}{Prop}}{}
\@ifundefined{lem}{\newtheorem{lem}[prop]{Lemma}}{}
\makeatother
\newcommand{\PP}{\mathbb P}
\title[Vanishing of Brauer groups of moduli stacks of stable curves]{Vanishing of Brauer groups of moduli stacks of stable curves}
\author{Sebastian Bartling, Kazuhiro Ito}
\date{Source: arXiv:2412.20435v1 (CC BY 4.0)}
\begin{document}
\maketitle

\noindent\emph{Source and licence.} Vanishing of Brauer groups of moduli stacks of stable curves. Version arXiv:2412.20435v1 (CC BY 4.0), distributed under the Creative Commons Attribution 4.0 International licence, as recorded in the arXiv licence metadata for this exact version and shown on the abstract page https://arxiv.org/abs/2412.20435v1. Full rights notice: licenses/arxiv-2412.20435-CC-BY-4.0.txt.\par\smallskip

\noindent\textit{Editorial scope.} Counted unit: the Lem whose source label is \texttt{Lemma:defining equation of ramified locus} in the article (source0.tex lines 1731--1737, source label \texttt{Lemma:defining equation of ramified locus}), delivered together with the article's own complete original proof of it (source0.tex lines 1739--1756, one \texttt{proof} environment of 18 lines). No mathematics is written, repaired or re-derived: statement and proof are the article's own text line for line. Required context retained uncounted: none. Every definition, display and label of those lines is retained unchanged. Omitted from this package and not redistributed: 1--1730, 1738--1738, 1757--1950. Nothing inside a retained excerpt is omitted or altered: every retained body is delivered exactly as the article's lines are written, so no omission receipt of any kind accompanies this package. Citations: the retained excerpts carry no citation command at all, so no bibliography entry of the article is transcribed and no reference is redistributed. Reference adaptation: none was needed and none was made; every reference inside the delivered unit resolves inside it, so no unresolved reference or citation is produced. Printed slips kept literal: no printed word, symbol, hypothesis or number of the counted statement or of its proof was altered. Layout and class: the article's own class is \texttt{amsart} with packages including amssymb, amsmath, amsthm, xcolor, hyperref, stmaryrd, xy, tikz-cd, cleveref, comment, relsize, mathbbol, amsfonts, color; the wrapper is the lane's pinned \texttt{amsart} editorial wrapper at 10pt on A4, which restates the theorem-like environments the retained text needs and adds the article's own macro definitions that the retained text needs (\texttt{\textbackslash PP}), with extra wrapper packages (bm, bbm, dsfont, xspace, cancel, mathtools). The delivered numbering is the unit's own. Assets: no retained span contains a figure environment, a graphics-inclusion command, a PDF-page inclusion, a table or a TikZ picture; no image file is copied into this package, the package declares no asset dependency and no image file is redistributed with it. Licence: version arXiv:2412.20435v1 of this article is distributed under the Creative Commons Attribution 4.0 International licence, as recorded in the arXiv licence metadata for this exact version and shown on the abstract page https://arxiv.org/abs/2412.20435v1. A full rights notice is delivered at licenses/arxiv-2412.20435-CC-BY-4.0.txt. Characters: every retained excerpt is pure ASCII, so the delivered engine, XeLaTeX with its default Latin Modern text font, reports no missing character.
\begin{lem}\label{Lemma:defining equation of ramified locus}
The curve
$C$ is defined by the equation
\[
g(X, Y, W):=f_1(X, Y, W)^2-4f_0(X, Y, W)f_2(X, Y, W)=0.
\]
\end{lem}

\begin{proof}
This is a well-known fact, but let us sketch the proof for the convenience of the reader.
We first note that the morphism
$\mathrm{Bl}_P (X) \to \PP^2_{F}$ maps the exceptional divisor isomorphically onto the line in $\PP^2_{F}$ defined by $f_0=0$.
The locus in $\PP^2_{F}$ where the geometric fibers of the projection $X - \{ P \} \to \PP^2_{F}$ are empty is the locus where $f_0=f_1=0$.
(The locus where $f_0=f_1=f_2=0$ is empty by the condition that $P \in X$ is not contained in any line in $X_{\overline{F}}$.)
The locus in $\PP^2_{F}$ where the geometric fibers of $X - \{ P \} \to \PP^2_{F}$ have exactly one point is the union of 
\[
\{ \, [x, y, z] \in \PP^2_{F} \, \vert \, f_0(x, y, z)=0, \, f_1(x, y, z) \neq 0 \, \}
\]
and
\[
\{ \, [x, y, z] \in \PP^2_{F} \, \vert \, f_0(x, y, z) \neq 0, \, g(x, y, z) = 0 \, \}.
\]
It follows that the ramification locus $C \subset \PP^2_{F}$ (or equivalently, the locus where the geometric fibers of
$\mathrm{Bl}_P (X) \to \PP^2_{F}$
have exactly one point) is defined by $g=0$.
\end{proof}

\end{document}
